\subsection{Theory of Constraints}\label{theory-of-constraints}

A constraint is any factor that prevents a system from achieving its
goal (Mabin and Balderstone, 2003). In manufacturing environments,
constraints can take the form of equipment with limited capacity or low
reliability, operations that produce high rates of defective parts, or
policies that limit the system's flexibility or efficiency (Gupta and
Boyd, 2008). Addressing the system's constraint often involves five
major steps (Watson \emph{et al.}, 2007): (i) identify the system's
constraint; (ii) exploit the constraint; (iii) subordinate everything
else to the constraint; (iv) elevate the constraint; and (v) if the
constraint is resolved, restart the process from step (i).

When demand \(d\) exceeds the planned capacity \(p\), identifying the
system's constraint (step 1) requires locating the resource
\(j \in \left\{ 0,\ 1,\ 2,\ \ldots,J \right\}\) with the highest load
\(l_{j}\), defined as \({argmax}_{j}l_{j}\) (\colorbox{yellow!40}{Reference}). In a
stochastic manufacturing environment producing multiple products
\(i \in \left\{ 0,\ 1,\ 2,\ \ldots,I \right\}\) in quantities \(m_{i}\),
utilizing \(J\) resources with expected cycle times
\(\mathbb{E}\left\lbrack \mathcal{T}_{ij}^{c} \right\rbrack\)\hspace{0pt},
considering each resource's expected overall effectiveness
\(\mathbb{E}\left\lbrack \mathcal{E}_{j} \right\rbrack\)\hspace{0pt}\hspace{0pt}
and change-over time
\(\mathbb{E}\left\lbrack \mathcal{T}_{j}^{CT} \right\rbrack\), the load
on a resource can be calculated as follows (\colorbox{yellow!40}{Reference}):

\begin{equation}
d = \left. \ \sum_{\forall i}^{}{m_{i}\max\left( \frac{\mathbb{E}\left\lbrack \mathcal{T}_{ij}^{c} \right\rbrack}{\mathbb{E}\left\lbrack \mathcal{E}_{j} \right\rbrack} + \frac{\mathbb{E}\left\lbrack \mathcal{T}_{j}^{CT} \right\rbrack}{m_{i}} \right)} \right|D_{i} \in D \land D = d;
\tag{8}
\end{equation}

\begin{equation}
\left. \ l_{j} = \frac{1}{p}\left\lbrack \left( \sum_{\forall i}^{}{m_{i}\frac{\mathbb{E}\left\lbrack \mathcal{T}_{ij}^{c} \right\rbrack}{\mathbb{E}\left\lbrack \mathcal{E}_{j} \right\rbrack}} \right)\mathbb{+ E}\left\lbrack \mathcal{T}_{j}^{CT} \right\rbrack \right\rbrack \right|d \geq p.
\tag{9}
\end{equation}

Once the constraint is identified, the next step is to determine the
most effective way to exploit it (step 2) and then adjust all other
operations accordingly (step 3) (Cox and Schleier, 2010). In production
scheduling, this is achieved by implementing DBR and BM (Telles \emph{et
al.}, 2020a). The constraint corresponds to the drum in DBR, which
dictates the overall system's performance (Darlington \emph{et al.},
2015). In TOC, manufacturing performance is primarily measured by
throughput \((\tau)\), which represents the money generated from sales
(Stefano \emph{et al.}, 2022). Throughput is calculated as the
difference between revenue \((\rho)\) and total variable costs \((c)\)
for each product \(i\), multiplied by its respective quantities produced
and sold. This can be expressed as:

\begin{equation}
\tau = \sum_{i = 1}^{I}{\left( \rho_{i} - c_{i} \right)m_{i}}.
\tag{10}
\end{equation}

Additionally, there are two other performance measures (\colorbox{yellow!40}{Reference}):
(i) inventory \(\mathcal{(i)}\), which represents the total money
invested in purchasing items intended for sale; and (ii) operating
expenses \((e)\), which reflect the money spent to convert inventory
into throughput.

Buffers are another class of element in DBR, typically categorized into
two types (when there are no assembly stations). The first is the
constraint buffer (CB), located before the constraint to protect its
flow and ensure it operates without shortages (\colorbox{yellow!40}{Reference}). The
second is the shipping buffer (SB), positioned at the end of the system
to protect against disruptions to the flow at any point in the process
(\colorbox{yellow!40}{Reference}). The initial size of these buffers
\(\left( b_{i}^{x} \right)\) for each product \(i\) is typically defined
by the quantity of semi-finished or finished products corresponding to
half the expected lead time from the start of the process \(j = 1\) to
the position \(j = n\), where the buffer
\(\mathcal{x} \in \left\{ CB,SB \right\}\) is located (\colorbox{yellow!40}{cf.,
Reference}). The lead time \(\mathcal{T}_{i}^{\mathcal{l}}\) is a random
variable resulting from the sum of all cycle times
\(\mathcal{T}_{ij}^{c}\) and waiting times \(\mathcal{T}_{ij}^{w}\)
belonging to a specific production route \(1 \rightarrow n\), and its
expected value is defined as:

\begin{equation}
\mathbb{E}\left\lbrack \mathcal{T}_{i}^{\mathcal{l}}\mathcal{|x} \right\rbrack\  = \ \sum_{\forall j \in (1 \rightarrow n)}^{}{\mathbb{E}\left\lbrack \mathcal{T}_{ij}^{c} \right\rbrack\mathbb{+ E}\left\lbrack \mathcal{T}_{ij}^{w} \right\rbrack}\left\{ \binom{n = {argmax}_{j}l_{j}\ }{n = N\ \ \ \ \ \ \ \ \ \ \ \ \ \ \ \ \ \ }\binom{if\mathcal{\ x =}CB}{if\mathcal{\ x =}SB} \right.\ .
\tag{11}
\end{equation}

According to Little's Law (\colorbox{yellow!40}{Reference}), the initial buffer size can
be estimated as follows:

\begin{equation}
b_{i}^{\mathcal{x}} \approx \min\left( \frac{\mathbb{E}\left\lbrack \mathcal{E}_{j} \right\rbrack}{\mathbb{E}\left\lbrack \mathcal{T}_{ij}^{c} \right\rbrack} \right)\left( \frac{\mathbb{E}\left\lbrack \mathcal{T}_{i}^{\mathcal{l}} \right\rbrack}{2} \right).
\tag{12}
\end{equation}

The initial buffer size is often dynamically adjusted to a `steady
state' reference level \({b'}_{i}^{x}\) using Buffer Management (BM)
policies (Schragenheim, 2010a). Within the TOC framework, the buffers
are typically divided into three distinct zones---red, yellow, and
green. Each zone represents one-third of the current reference level,
with the red zone indicating the lowest inventory level and the green
zone the highest (Telles \emph{et al.}, 2022b). The reference leve is
continuous evaluated over a replenishment time \(\mathcal{T}_{i}^{R}\),
that is difined as the sum of order planning interval
\(\mathrm{\Delta}\varphi\) and the lead time
\(\mathcal{T}_{i}^{\mathcal{l}}\) (Schragenheim, 2010b). If the actual
buffer level, \(H_{i}^{x}\), remains consistently within the green zone
for a duration of at least three replenishment times
\(3\mathcal{T}_{i}^{R}\), the reference level is reduced by 15\% for the
subsequent order planning interval, \(\mathrm{\Delta}\varphi + 1\). The
reduction logic is defined as:

\begin{equation}
{b'}_{i,\Delta\varphi + 1}^{x} = 0.85b_{i,\Delta\varphi}^{x}\quad\text{if}\quad H_{i,t}^{x} \geq \frac{2}{3}b_{i,\Delta\varphi}^{x}\ |\ \forall t \in \left\lbrack 0,\ 3\mathcal{T}_{i}^{R} \right\rbrack
\tag{13}
\end{equation}

Conversely, the buffer must be expanded if the cumulative penetrations
into the red zone, \(k_{i}^{x}\), achieve or exceed the total size of
the red zone during a single replenishment time \(\mathcal{T}_{i}^{R}\)
(Schragenheim, 2010b). The reference level is increased by 20\% for the
next order planning interval \(\Delta\varphi\). The expansion logic is
defined as:

\begin{equation}
{b'}_{i,\Delta\varphi + 1}^{x} = 1.2b_{i,\Delta\varphi}^{x}\quad\text{if}\quad\sum_{t = 0}^{\mathcal{T}_{i}^{R}}k_{i,t}^{x} \geq \frac{1}{3}b_{i,\Delta\varphi}^{x}
\tag{14}
\end{equation}

The rationale for these adjustments is to protect the production flow
while maintaining the minimum inventory \(\mathcal{i}\).

Similar to buffers, production scheduling is another approach to
exploiting the constraint (step 2) (\colorbox{yellow!40}{Reference}). In TOC, the goal of
scheduling is to maximize throughput during the replenishment interval
(\colorbox{yellow!40}{Reference}). Therefore, the sequence of products \(i\) to be
produced is defined in descending order of throughput, moderated by the
priority level of finished products in the SB, i.e.,
\({\left( \frac{\mathcal{b'}_{i}^{SB}}{H_{i}^{SB}} \right)\tau}_{i}^{(1)} \geq \left( \frac{\mathcal{b'}_{i}^{SB}}{H_{i}^{SB}} \right)\tau_{i}^{(2)} \geq \ldots \geq \left( \frac{\mathcal{b'}_{i}^{SB}}{H_{i}^{SB}} \right)\tau_{i}^{(\mathfrak{n})}\),
based on the following equation:

\begin{equation}
\left. \ \underset{m_{i}}{argmax\ }\tau \right|\frac{\mathrm{\Delta}d}{\mathrm{\Delta}\varphi} \leq \frac{\mathrm{\Delta}p}{\mathrm{\Delta}\varphi}.
\tag{15}
\end{equation}

The rope is the final element of DBR, and its function is to subordinate
the rest of the system to the constraint (step 3) (\colorbox{yellow!40}{Reference}).
Often, the demand is greater than the capacity during the replenishment
interval, which could lead to the creation and release of orders that
the system cannot fully execute, ultimately increasing \(\mathcal{i}\)
(\colorbox{yellow!40}{Reference}). Therefore, the rope's role is to release production
orders based on the consumption of the CB, as outlined in
\colorbox{yellow!40}{\hyperref[_Ref187516903]{Figure 3}}.

The quantity of each product \(i\) to be produced during the
replenishment interval \(M_{i}^{\mathrm{\Delta}\varphi}\) is typically a
function of the SB level \(H_{i}^{SB}\) and its reference
\(\mathcal{b'}_{i}^{SB}\) and the work-in-process \(W_{i}^{SB}\) along a
complete production route \(1 \rightarrow N\), as described by the
following equations (Cox III and Schleier Jr., 2010):

\begin{equation}
W_{i}^{SB} = \sum_{\forall j \in (1 \rightarrow N)}^{}w_{ij};
\tag{16}
\end{equation}

\begin{equation}
M_{i}^{\mathrm{\Delta}\varphi} = \left( \mathcal{b'}_{i}^{SB} - H_{i}^{SB} \right) - W_{i}^{SB};
\tag{17}
\end{equation}

where \(w_{ij}\) represents the quantity of product \(i\) being produced
at the resource \(j\). This quantity is stored in a `virtual buffer'
known as the order buffer (OB), which typically has an initial infinite
size \(\left( \mathcal{b}_{i}^{OB} = \infty \right)\) and a reference
level based on the replenishment interval, i.e.,
\({\mathcal{b}'}_{i}^{OB} = M_{i}^{\mathrm{\Delta}\varphi}\).
Additionally, the quantity of each product \(i\) to be released for
production within a given time interval \(M_{i}^{\mathrm{\Delta}t}\) is
a function of the CB level \(H_{i}^{CB}\), its reference
\(\mathcal{b'}_{i}^{CB}\), and the work-in-process \(W_{i}^{CB}\) along
a partial production route \(1 \rightarrow {argmax}_{j}l_{j}\) that ends
at the constraint. This relationship is described as follows:

\begin{equation}
W_{i}^{CB} = \sum_{\forall j \in \left( 1 \rightarrow {argmax}_{j}l_{j} \right)}^{}w_{ij};
\tag{18}
\end{equation}

\begin{equation}
M_{i}^{\mathrm{\Delta}t} = \left( \mathcal{b'}_{i}^{CB} - H_{i}^{CB} \right) - W_{i}^{CB}.
\tag{19}
\end{equation}

The number of orders released reduces the current level of the OB,
denoted as \(H_{i}^{OB}\).

Step 4 may involve increasing shifts, hiring additional labor, or
purchasing new equipment to increase the capacity of the constraint
(Goldratt, 2006b). Step 5 continues the cycle of continuous improvement
within TOC (Goldratt and Cox, 1984b). Both of these steps are beyond the
scope of the current work.
